% Part: first-order-logic
% Chapter: tableaux
% Section: provability-consistency

\documentclass[../../../include/open-logic-section]{subfiles}

\begin{document}

\iftag{FOL}
      {\olfileid{fol}{tab}{prv}}
      {\olfileid{pl}{tab}{prv}}

\olsection{\usetoken{S}{derivability} ও সঙ্গতি}

এখন !!{derivability}-সম্পর্কের কয়েকটি ধর্ম প্রতিষ্ঠা করব। এগুলি
স্বতন্ত্রভাবেও গুরুত্বপূর্ণ; আবার প্রত্যেকটি পূর্ণতা উপপাদ্যের
প্রমাণে কাজে লাগবে।

\begin{prop}\ollabel{prop:provability-contr}
  $\Gamma \Proves !A$ এবং $\Gamma \cup \{!A\}$ অসঙ্গত হলে
  $\Gamma$ অসঙ্গত।
\end{prop}

\begin{proof}
$\Gamma_0=\{!B_1,\dots,!B_n\}$ এবং
$\Gamma_1=\{!C_1,\dots,!C_m\}\subseteq\Gamma$ এমন সসীম সমষ্টি আছে,
যাদের জন্য
  \begin{align*}
    \{\sFmla{\False}{!A}, &
    \sFmla{\True}{!B_1}, \dots, \sFmla{\True}{!B_n}\} \\
    \{\sFmla{\True}{!A}, &
    \sFmla{\True}{!C_1}, \dots, \sFmla{\True}{!C_m}\}
  \end{align*}
-এর বদ্ধ !!{tableau}s আছে। $!A$-তে \Cut{} বিধি ব্যবহার করে এগুলিকে
একটি বদ্ধ !!{tableau}-তে একত্র করা যায়, যা দেখায়
$\Gamma_0\cup\Gamma_1$ অসঙ্গত। যেহেতু
$\Gamma_0\subseteq\Gamma$ এবং $\Gamma_1\subseteq\Gamma$, তাই
$\Gamma_0\cup\Gamma_1\subseteq\Gamma$; অতএব $\Gamma$ অসঙ্গত।
% BN-SRC-050 TARGET-CORRECTION-BEGIN
\par\smallskip\noindent\textnormal{\small\textbf{উৎস-সংশোধন (BN-SRC-050):} হিমায়িত ইংরেজি গদ্যে $\Gamma_1$-কে $\{!C_1,\dots,!C_n\}$ বলা হলেও একই প্রদর্শন ও পরের ব্যবহার শেষ সূচক~$m$ দেয়। $\Gamma_0$-র $n$টি এবং $\Gamma_1$-এর $m$টি বাক্য আলাদা রাখতে বাংলা লক্ষ্যে $\{!C_1,\dots,!C_m\}$ পুনরুদ্ধার করা হয়েছে।} % BN-SRC-050 TARGET-CORRECTION-END
\end{proof}

\begin{prop}
\ollabel{prop:prov-incons}
$\Gamma \Proves !A$ হবে তখনই এবং কেবল তখনই, যখন
$\Gamma \cup \{\lnot !A\}$ অসঙ্গত।
\end{prop}

\begin{proof}
প্রথমে ধরি $\Gamma \Proves !A$; অর্থাৎ নিচের সমষ্টিটির জন্য একটি
বদ্ধ !!{tableau} আছে:
\[
\{\sFmla{\False}{!A},
\sFmla{\True}{!B_1}, \dots, \sFmla{\True}{!B_n}\}
\]
$\TRule{\True}{\lnot}$ বিধি ব্যবহার করে একে নিচের সমষ্টিটির জন্য
একটি বদ্ধ !!{tableau}-তে বদলানো যায়:
\[
\{\sFmla{\True}{\lnot !A},
\sFmla{\True}{!B_1}, \dots, \sFmla{\True}{!B_n}\}.
\]

অন্যদিকে, শেষের সমষ্টিটির জন্য বদ্ধ !!{tableau} থাকলে সেটিকে প্রথম
সমষ্টিটির বদ্ধ !!{tableau}-তে বদলানো যায়। প্রথম অনুমিতি
$\sFmla{\True}{\lnot !A}$-তে \TRule{\True}{\lnot} প্রয়োগের ফলে
পাওয়া প্রতিটি সূত্র ও সেই অনুমিতি সরিয়ে দিই এবং অনুমিতি
$\sFmla{\False}{!A}$ যোগ করি। কোনো শাখা আগে
$\sFmla{\True}{\lnot !A}$-তে \TRule{\True}{\lnot}-এর প্রয়োগের সিদ্ধান্ত, অর্থাৎ
$\sFmla{\False}{!A}$ থাকার কারণে বদ্ধ হলে, নতুন !!{tableau}-এর
সংশ্লিষ্ট শাখাটিও বদ্ধ। পুরোনো ট্যাবলোর কোনো শাখা যদি অনুমিতি
$\sFmla{\True}{\lnot !A}$ ও $\sFmla{\False}{\lnot !A}$ উভয়ই থাকার
কারণে বদ্ধ হয়ে থাকে, তবে $\sFmla{\False}{\lnot !A}$-তে
$\TRule{\False}{\lnot}$ প্রয়োগ করে $\sFmla{\True}{!A}$ পাওয়া যায়
এবং শাখাটি বদ্ধ করা যায়। কারণ অনুমিতি হিসেবে
$\sFmla{\False}{!A}$ আগেই যোগ করেছি।
\end{proof}

\begin{prob}
$\Gamma \Proves \lnot !A$ হবে তখনই এবং কেবল তখনই, যখন
$\Gamma \cup \{!A\}$ অসঙ্গত—প্রমাণ করো।
\end{prob}

\begin{prop}\ollabel{prop:explicit-inc}
  $\Gamma \Proves !A$ এবং $\lnot !A \in \Gamma$ হলে $\Gamma$ অসঙ্গত।
\end{prop}

\begin{proof}
  ধরি $\Gamma \Proves !A$ এবং $\lnot !A \in \Gamma$। তখন
  $!B_1$, \dots, $!B_n\in\Gamma$ এমন বাক্য আছে যেন
  \[
  \{\sFmla{\False}{!A}, \sFmla{\True}{!B_1}, \dots, \sFmla{\True}{!B_n}\}
  \]
-এর একটি বদ্ধ ট্যাবলো আছে। অনুমিতি \sFmla{\False}{!A}-এর বদলে
  \sFmla{\True}{\lnot !A} বসাই এবং অনুমিতিগুলির পরে
  \sFmla{\True}{\lnot !A}-তে \TRule{\True}{\lnot} প্রয়োগের সিদ্ধান্ত
  যোগ করি। পুরোনো !!{tableau}-তে পংক্তি~$1$ উল্লেখ করে সমর্থিত
  প্রত্যেক !!{sentence} এখন পংক্তি~$n+1$ উল্লেখ করে সমর্থিত। অতএব
  পুরোনো !!{tableau}-টি বদ্ধ হলে নতুনটিও বদ্ধ। নতুন !!{tableau}-টি দেখায়
  $\Gamma$ অসঙ্গত, কারণ সব অনুমিতিই~$\Gamma$-তে আছে।
% BN-SRC-051 TARGET-CORRECTION-BEGIN
\par\smallskip\noindent\textnormal{\small\textbf{উৎস-সংশোধন (BN-SRC-051):} হিমায়িত ইংরেজি গদ্যে $\TRule{\True}{\lnot}$ বিধির সিদ্ধান্তটিকে ভুল করে $\sFmla{\False}{!A}$-তে বিধিটি প্রয়োগের ফল বলা হয়েছে। একই বাক্যে যে প্রতিস্থাপিত অনুমিতি দেওয়া আছে এবং বিধির সংজ্ঞা অনুসারে বাংলা লক্ষ্যে প্রয়োগের পূর্বধারণা $\sFmla{\True}{\lnot !A}$ পুনরুদ্ধার করা হয়েছে।} % BN-SRC-051 TARGET-CORRECTION-END
\end{proof}

\begin{prop}\ollabel{prop:provability-exhaustive}
  $\Gamma \cup \{!A\}$ এবং $\Gamma \cup \{\lnot !A\}$ উভয়ই
  অসঙ্গত হলে $\Gamma$ অসঙ্গত।
\end{prop}

\begin{proof}
  যদি $!B_1$, \dots, $!B_n\in\Gamma$ এবং
  $!C_1$, \dots, $!C_m\in\Gamma$ এমন বাক্য থাকে যেন
  \begin{align*}
    \{\sFmla{\True}{!A}, &
    \sFmla{\True}{!B_1}, \dots, \sFmla{\True}{!B_n}\} \text{ এবং}\\
    \{\sFmla{\True}{\lnot !A}, &
    \sFmla{\True}{!C_1}, \dots, \sFmla{\True}{!C_m}\}
  \end{align*}
উভয়েরই বদ্ধ !!{tableau}s থাকে, তবে একটি সম্মিলিত !!{tableau}
  তৈরি করে দেখানো যায় যে $\Gamma$ অসঙ্গত। অনুমিতি হিসেবে
  $\sFmla{\True}{!B_1}$, \dots, $\sFmla{\True}{!B_n}$-এর সঙ্গে
  $\sFmla{\True}{!C_1}$, \dots, $\sFmla{\True}{!C_m}$ নিই এবং তারপর
  \Cut{} বিধি প্রয়োগ করি। এতে দুটি শাখা হয়—একটি
  $\sFmla{\True}{!A}$ দিয়ে, অন্যটি $\sFmla{\False}{!A}$ দিয়ে শুরু।

  বাঁ দিকে প্রথম !!{tableau}-টির অনুমিতির নিচের অংশ যোগ করি। এখানে
  প্রতিটি বিধি-প্রয়োগ এখনও সঠিক, কারণ
  $\sFmla{\True}{!A}$-সহ প্রথম !!{tableau}-টির প্রতিটি অনুমিতি পাওয়া
  যাচ্ছে। ফলে $\sFmla{\True}{!A}$-এর নিচের প্রতিটি শাখা বদ্ধ হয়।
% BN-SRC-052 TARGET-CORRECTION-BEGIN
\par\smallskip\noindent\textnormal{\small\textbf{উৎস-সংশোধন (BN-SRC-052):} হিমায়িত ইংরেজি গদ্যে “left left side” মুদ্রণপ্রমাদ আছে। বৃক্ষের দ্বিবিভাজন অনুসারে বাংলা লক্ষ্যে একবারই “বাঁ দিক” লেখা হয়েছে।} % BN-SRC-052 TARGET-CORRECTION-END

  ডান দিকে দ্বিতীয় !!{tableau}-টির অনুমিতির নিচের অংশ যোগ করি এবং
  $\sFmla{\True}{\lnot !A}$-তে~$\TRule{\True}{\lnot}$ প্রয়োগের
  ফলগুলি সরিয়ে দিই। $\sFmla{\True}{\lnot !A}$-তে
  $\TRule{\True}{\lnot}$ প্রয়োগের সিদ্ধান্ত হলো
  $\sFmla{\False}{!A}$; তবু সেটি পাওয়া যাচ্ছে, কারণ সম্মিলিত
  !!{tableau}-এর ডান দিকে সেটিই \Cut{} বিধির সিদ্ধান্ত।

  দ্বিতীয় ট্যাবলোর কোনো শাখা যদি অনুমিতি
  $\sFmla{\True}{\lnot !A}$ (যা সম্মিলিত !!{tableau}-তে আর অনুমিতি
  নয়) এবং $\sFmla{\False}{\lnot !A}$ উভয়ই থাকার কারণে বদ্ধ হয়ে
  থাকে, তবে $\sFmla{\False}{\lnot !A}$-তে
  $\TRule{\False}{\lnot}$ প্রয়োগ করে $\sFmla{\True}{!A}$ পাই। এখন
  সম্মিলিত !!{tableau}-এর সংশ্লিষ্ট শাখাটিও বদ্ধ, কারণ তাতে
  \Cut{} বিধির ডানদিকের সিদ্ধান্ত~$\sFmla{\False}{!A}$ আছে। দ্বিতীয়
  !!{tableau}-এর কোনো শাখা অন্য কোনো কারণে বদ্ধ হলে সম্মিলিত
  !!{tableau}-এর সংশ্লিষ্ট শাখাটিও বদ্ধ; কারণ পুরোনো দ্বিতীয়
  !!{tableau}-এর শাখায় $\sFmla{\True}{\lnot !A}$ ছাড়া থাকা অন্য
  সব !!{signed formula}s সম্মিলিত !!{tableau}-এর সংশ্লিষ্ট শাখাতেও
  থাকে।
  \end{proof}

\end{document}
